% !TeX program = tectonic
\documentclass[11pt,a4paper]{article}
\usepackage{fontspec}
\usepackage[english]{babel}
\babelfont{rm}[Language=Default]{Liberation Serif}
\babelfont[Language=Default]{sf}{Carlito}
\babelfont[Language=Default]{tt}{DejaVu Sans Mono}
\usepackage{amsmath,amssymb,amsthm}
\usepackage{geometry}
\geometry{a4paper, top=2.4cm, bottom=2.4cm, left=2.4cm, right=2.4cm}
\usepackage{booktabs,tabularx,enumitem,xcolor,tcolorbox}
\tcbuselibrary{breakable,skins}
\definecolor{linkblue}{HTML}{1F4E79}
\definecolor{statgreen}{HTML}{2F6B3A}
\definecolor{goldacc}{HTML}{8B7E5A}
\usepackage[colorlinks=true,linkcolor=linkblue,urlcolor=linkblue,unicode=true]{hyperref}
\theoremstyle{plain}
\newtheorem{theorem}{Theorem}
\newtheorem{lemma}[theorem]{Lemma}
\newtheorem{proposition}[theorem]{Proposition}
\newtheorem{corollary}[theorem]{Corollary}
\theoremstyle{definition}
\newtheorem{definition}[theorem]{Definition}
\theoremstyle{remark}
\newtheorem{remark}[theorem]{Remark}
\newtcolorbox{statusbox}[1][]{colback=statgreen!5,colframe=statgreen!75,
  title={\textbf{Summary of results:} #1},fonttitle=\small,boxrule=0.7pt,arc=2pt,breakable}
\newtcolorbox{databox}[1][]{colback=goldacc!6,colframe=goldacc!80,
  title={\textbf{Protocol data:} #1},fonttitle=\small,boxrule=0.7pt,arc=2pt,breakable}
\newtcolorbox{verifybox}[1][]{colback=linkblue!5,colframe=linkblue!80,
  title={\textbf{Verification:} #1},fonttitle=\small,boxrule=0.7pt,arc=2pt,breakable}
\widowpenalty=10000
\clubpenalty=10000
\setlength{\parskip}{0.35em}
\newcommand{\CC}{\mathbb{C}}
\newcommand{\RR}{\mathbb{R}}
\newcommand{\ZZ}{\mathbb{Z}}
\newcommand{\QQ}{\mathbb{Q}}
\newcommand{\Ga}{\Gamma}
\newcommand{\im}{\operatorname{Im}}
\newcommand{\re}{\operatorname{Re}}
\newcommand{\disc}{\operatorname{disc}}
\begin{document}
\begin{center}
{\LARGE\bfseries The Conductor Census: Two Schemes, One Genus}\\[0.4em]
{\large Exact character histograms across the four rungs $N=7/9/15/30$}\\[0.6em]
{\normalsize Isaev Iskhak Khamzatovich}\\[0.2em]
{\small The ``Dynamic Principle'' program --- the \texttt{hodge-laboratory} repository}\\[0.15em]
{\small Standalone theorem monograph · Monograph 20/20}
\end{center}
\vspace{0.4em}
{\color{linkblue}\hrule height 0.8pt}
\vspace{0.8em}

\section{Setting}

While monograph 1/20 derived the genus of the Fermat curve
topologically (Riemann--Hurwitz), here the genus is obtained
\emph{arithmetically} --- as the total of the conductor census of
the eigencharacters. The census is the arithmetic skeleton of the
ladder: it determines the multiplicities $h_d$ with which the
conductors $d\mid N$ enter every Coleman--Selberg product of the
program (certificates G and H), and it works identically on all
four rungs $N=7,9,15,30$. Protocol V1 executes \emph{two
independent census schemes} --- the direct one and the
M\"obius one --- and demands their bit-for-bit agreement; this
monograph proves that both schemes count the same thing and
collects the exact histograms of every rung.

\begin{definition}[eigencharacter and conductor]\label{th:def}
An eigencharacter of level $N$ is a pair $(a,b)$ of integers with
$1\le a$, $1\le b$, $a+b\le N-1$; it defines a holomorphic
differential $\omega_{a,b}=x^ay^b\,dx/y^b$ on the Fermat curve
$X_N\colon x^N+y^N=1$. The conductor of a character is
$d=N/\gcd(N,a,b)$ --- the smallest level through which the
character factors. The census is the histogram $h_d$ of the
conductors; $c(k)=(k-1)(k-2)/2$ is the total number of
eigencharacters of level $k$.
\end{definition}

\begin{theorem}[the triangle]\label{th:tri}
$c(N)=(N-1)(N-2)/2$, and the census total equals the genus:
$\sum_{d\mid N}h_d=c(N)=(N-1)(N-2)/2=g$.
\end{theorem}

\begin{proof}
The number of pairs $1\le a,b$ with $a+b\le N-1$ is the count of
lattice points of the integer triangle: for $a=1,\dots,N-2$ there
are exactly $N-1-a$ values of $b$, giving
$\sum_{a=1}^{N-2}(N-1-a)=1+2+\dots+(N-2)=(N-1)(N-2)/2$.
Each pair yields a basis holomorphic differential (classical for
the Fermat curve, monograph 1/20), so $c(N)=g$ --- the
arithmetic side of the genus formula, independent of
Riemann--Hurwitz. Since every character has exactly one conductor
(lemma~\ref{th:lift} below shows the conductors divide $N$), the
census total is $c(N)$.
\end{proof}

\begin{theorem}[the lifting lemma and M\"obius inversion]\label{th:lift}
For every $d\mid N$ the map
$(a',b')\mapsto\big(\tfrac{N}{d}a',\tfrac{N}{d}b'\big)$ is a
bijection between the eigencharacters of level $d$ and the
eigencharacters of level $N$ whose conductor divides $d$.
Consequently,
\[
  \sum_{d'\mid d}h_{d'}=c(d),
  \qquad\text{and by M\"obius inversion}\qquad
  h_d=\sum_{m\mid d}\mu(m)\,c(d/m).
\]
In particular $h_d$ does not depend on the ambient level $N$
(only $d\mid N$ matters): $h_d(N)=h_d(d)$ --- \emph{the
self-similarity of the ladder}: every conductor carries a
complete smaller stand.
\end{theorem}

\begin{proof}
\emph{Well-definedness.} Let $(a',b')$ be a character of level
$d$. Put $a=\tfrac{N}{d}a'$, $b=\tfrac{N}{d}b'$: then
$1\le a,b$ and $a+b=\tfrac{N}{d}(a'+b')\le\tfrac{N}{d}(d-1)
=N-\tfrac{N}{d}\le N-1$, so $(a,b)$ is a character of level $N$.
Its conductor: $\gcd(N,a,b)=\tfrac{N}{d}\gcd(d,a',b')$, hence
$d(a,b)=N/\gcd(N,a,b)=d/\gcd(d,a',b')=d(a',b')\mid d$.

\emph{Surjectivity with a section.} If a character $(a,b)$ of
level $N$ has conductor $d'\mid d$, then
$\gcd(N,a,b)=\tfrac{N}{d'}$ is a multiple of $\tfrac{N}{d}$,
whence $\tfrac{N}{d}\mid a,b$ and the pair $a'=\tfrac{d}{N}a$,
$b'=\tfrac{d}{N}b$ is integral; the bounds are preserved by the
reverse of the well-definedness argument, and $(a',b')$ is a
character of level $d$ with conductor $d'$ that restores $(a,b)$.
Injectivity is obvious. Therefore
$\sum_{d'\mid d}h_{d'}(N)=c(d)$ --- the total number of
characters of level $d$.

\emph{Inversion.} The divisor lattice of $d$ is a finite
distributive lattice, and its M\"obius function is the classical
$\mu$; M\"obius inversion on the divisor lattice yields the
second formula. The right-hand side depends only on $d$, so $h_d$
is the same at every level $N$ divisible by $d$: at the level $d$
the formula produces $h_d(d)$. The self-similarity is proved.
\end{proof}

\begin{theorem}[the censuses of the four rungs]\label{th:four}
The exact conductor histograms:
\[
\begin{array}{lll}
N=7: & \{7\colon 15\}, & 15=15;\\[0.15em]
N=9: & \{3\colon 1,\ 9\colon 27\}, & 28=1+27;\\[0.15em]
N=15: & \{3\colon 1,\ 5\colon 6,\ 15\colon 84\}, & 91=1+6+84;\\[0.15em]
N=30: & \{3\colon 1,\ 5\colon 6,\ 6\colon 9,\ 10\colon 30,\ 15\colon 84,\ 30\colon 276\}, & 406=1+6+9+30+84+276.
\end{array}
\]
Each row is verified by both schemes, including the explicit
values $h_6=c(6)-c(3)-c(2)+c(1)=10-1-0+0=9$,
$h_{10}=c(10)-c(5)-c(2)+c(1)=36-6=30$ and
$h_{30}=c(30)-c(15)-c(10)-c(6)+c(5)+c(3)+c(2)-c(1)
=406-91-36-10+6+1=276$.
\end{theorem}

\begin{proof}
\emph{The direct scheme.} $N=7$ is prime: $\gcd(7,a,b)=1$ for
every character of the triangle --- the single conductor $7$,
$h_7=c(7)=15$. $N=9$: conductor $3$ means $\gcd(9,a,b)=3$, i.e.
$a,b\in\{3,6\}$, and $a+b\le8$ leaves $(3,3)$: $h_3=1$; the rest
$h_9=28-1=27$. $N=15$: conductor $3$ --- $\gcd(15,a,b)=5$, i.e.
$a,b\in\{5,10\}$ with $a+b\le14$ --- only $(5,5)$: $h_3=1$;
conductor $5$ --- $\gcd=3$, i.e. $a,b\in\{3,6,9,12\}$ with
$a+b\le14$ --- the six pairs $(3,3),(3,6),(6,3),(3,9),(9,3),(6,6)$:
$h_5=6$; the rest $h_{15}=91-1-6=84$. $N=30$ --- conductor $6$:
$\gcd(30,a,b)=5$, so $a,b$ are multiples of $5$ with $a+b\le29$
and $\gcd$ exactly $5$ (not $10$, not $15$) --- nine pairs; the
M\"obius scheme gives the same $h_6=9$; likewise $h_{10}=30$,
$h_{15}=84$ (self-similarity: $h_{15}(30)=h_{15}(15)$) and
$h_{30}=276$.

\emph{The M\"obius scheme.} The formula of
theorem~\ref{th:lift} with the explicit values $c$: $c(2)=0$,
$c(3)=1$, $c(5)=6$, $c(6)=10$, $c(9)=28$, $c(10)=36$,
$c(15)=91$, $c(30)=406$; the computations are displayed in the
statement --- every conductor matches the direct scheme
pairwise.
\end{proof}

\begin{theorem}[two schemes, one total]\label{th:two}
The direct scheme (enumerating the pairs with
$d=N/\gcd(N,a,b)$) and the M\"obius scheme
(inclusion--exclusion over the divisors) produce
\emph{bit-for-bit identical} histograms at every level
$N=7,9,15,30$ --- an integer identity, a deterministic check
without thresholds.
\end{theorem}

\begin{proof}
By theorem~\ref{th:lift} both schemes compute the same function
$h_d$: the direct one by definition (every pair lands in exactly
one conductor class), the M\"obius one as the unique solution of
the convolution equation $\sum_{d'\mid d}h_{d'}=c(d)$ on the
divisor lattice. The formal coincidence leaves no room for
discrepancies; the machine execution (protocol V1) is a
fingerprint against implementation errors: two independent code
paths must produce identical dictionaries $\{d\colon h_d\}$,
which is exactly what the bit-wise comparison asserts.
\end{proof}

\begin{remark}
The histograms compress quadratically with the arithmetic of the
level: the prime $7$ has one conductor; $9=3^2$ has two (the
``hermit'' $(3,3)$ appears --- the trace of level $3$);
$15=3\cdot5$ has three; $30=2\cdot3\cdot5$ has six --- every
divisor of the level except $1$ and $2$ (the triangle of level
$d$ is empty only for $d\le2$). The self-similarity
$h_d(N)=h_d(d)$ turns the ladder into a fractal family of
stands: a larger rung contains all smaller ones as
conductor layers, and every summand of
$406=1+6+9+30+84+276$ is itself the census of a smaller stand
(its proper conductors subtracted). It is the multiplicities
$h_d$ that enter the exponents of the Coleman--Selberg product of
certificate H --- the census closes the $\Ga$-architecture of the
program.
\end{remark}

\begin{databox}[The census --- a summary of the four rungs]
$N=7$: genus $15$, $\{7\colon15\}$, a single conductor (prime
level). $N=9$: genus $28$, $\{3\colon1,\,9\colon27\}$, the
diagonal character $(3,3)$. $N=15$: genus $91$,
$\{3\colon1,\,5\colon6,\,15\colon84\}$. $N=30$: genus $406$,
$\{3\colon1,\,5\colon6,\,6\colon9,\,10\colon30,\,15\colon84,\,30\colon276\}$.
All rows confirmed by two schemes (V1); the totals equal
$(N-1)(N-2)/2$ exactly; the self-similarity $h_d(N)=h_d(d)$
(e.g. $h_{15}(30)=84=h_{15}(15)$). Reference:
\texttt{baseline\_v1\_v9.json}, block \texttt{V1\_census}.
\end{databox}

\begin{statusbox}[summary]
The arithmetic of the genus is placed on its own foundation: the
character triangle, the lifting lemma with the conductor
bijection, M\"obius inversion, the self-similarity
$h_d(N)=h_d(d)$, the exact histograms of all four rungs and the
bit-for-bit agreement of the two schemes. The genus is
\emph{counted}, not postulated; the multiplicities $h_d$ are the
input of the $\Ga$-normalization of certificates G/H.
\end{statusbox}

\begin{verifybox}[how to reproduce]
\texttt{python3 laboratory.py} $\to$ protocol V1 (the census,
two schemes, levels $7/9/15/30$); \texttt{--check-baseline}
cross-checks the histograms and the genera against the
reference; Lean: the \texttt{genus\_fermat} block; batch mode:
the \texttt{census} item of a scenario.
\end{verifybox}

\vfill
{\color{linkblue}\hrule height 0.6pt}
\vspace{0.5em}
{\small\color{gray}
License: individual exclusive license of Isaev Iskhak Khamzatovich.
All rights to the computations, formulas and texts belong to the author.
Verification: \texttt{python3 laboratory.py} (protocol V1, \texttt{--check-baseline}).
}
\end{document}
